\chapter{Design Space of Graph Data Management Systems}
\label{chap:design_space}

In the previous chapter~\autoref{chapter:background}, we discussed the different models used to represent graphs -- property graphs and how they can be used to materialize a structural graph. 
We also discussed the different types of workloads that can be executed on these graphs and the access patterns they exhibit.
Based on these access patterns, we identified requirements for a data structure that can be used to store graphs and discussed the different data structures that are used by graph data management systems.

In this chapter, we present a comprehensive overview of the design space of graph data management systems. We begin by discussing the different data layouts used by these systems, both in-memory and persistent. We then delve into the different workloads that these systems are designed to handle, and the different types of graph databases that are available. We also discuss the different computational models that are used by these systems. Finally, we present a detailed discussion of the different graph analytics systems that are available, and the different optimizations that they use to reduce the communication overhead.

